\chapter{Realizacja logiki logistycznej i algorytmów}

W niniejszym rozdziale przedstawiono szczegóły implementacyjne kluczowych algorytmów sterujących przepływem przesyłek w systemie FasterPost. Logika została podzielona na dwie główne warstwy: logistykę krajową (transport między magazynami) oraz logistykę lokalną (dostawy ostatniej mili do Paczkomatów).

\section{Logistyka krajowa: Implementacja problemu CVRP}

Logistyka krajowa odpowiada za transport przesyłek pomiędzy magazynami głównymi (Hub-to-Hub). Problem ten został zamodelowany jako wariant problemu trasowania pojazdów z ograniczeniami pojemności (CVRP - Capacitated Vehicle Routing Problem).

Implementacja znajduje się w klasie \texttt{RoutingService} w module \texttt{logistics}. Algorytm uwzględnia dwa kluczowe ograniczenia:
\begin{itemize}
    \item \textbf{Pojemność pojazdu} (\texttt{vehicle\_capacity}): Maksymalna liczba przesyłek, jaką może zabrać jeden kurier (domyślnie 50).
    \item \textbf{Czas pracy} (\texttt{MAX\_WORK\_DAY\_MINUTES}): Maksymalny czas trwania trasy, ustalony na 720 minut (12 godzin).
\end{itemize}

\subsection{Heurystyka Najbliższego Sąsiada (Nearest Neighbor)}

Ze względu na złożoność obliczeniową problemu CVRP (NP-trudny), zastosowano podejście heurystyczne typu zachłannego (ang. greedy algorithm). Metoda \texttt{\_build\_optimized\_tour} iteracyjnie wybiera najbliższy magazyn docelowy, do którego kurier może dotrzeć, nie przekraczając limitu czasu pracy.

\begin{lstlisting}[language=Python, caption=Fragment algorytmu wyboru kolejnego celu w logistyce krajowej]
while remaining_dests:
    best_next = None
    best_dist = float('inf')
    
    # Wybor najblizszego magazynu (Nearest Neighbor)
    for candidate_wh in remaining_dests:
        d = self.distance_service.get_distance(
            str(current_node.id), str(candidate_wh.id)
        )
        if not math.isinf(d) and d < best_dist:
            best_dist = d
            best_next = candidate_wh
\end{lstlisting}

System wykorzystuje również algorytm przeszukiwania wszerz (BFS) zaimplementowany w metodzie \texttt{\_find\_shortest\_graph\_path}, aby znaleźć drogę w grafie połączeń między magazynami, jeśli nie istnieje bezpośrednie połączenie.

\section{Logistyka lokalna: Ostatnia mila i TSP}

Logistyka lokalna (\textit{ang. last mile delivery}) obsługiwana jest przez klasę \texttt{LocalRoutingService}. Proces ten składa się z trzech etapów: grupowania przesyłek według stref, alokacji skrytek oraz wyznaczania optymalnej trasy kuriera.

\subsection{Klasteryzacja i podział na strefy}

System operuje na pojęciu stref (\texttt{Zone}), które reprezentują terytoria przypisane do konkretnego magazynu. Choć definicja stref jest statyczna w momencie planowania trasy, ich kształt wynika z wcześniejszej analizy gęstości Paczkomatów (koncepcyjnie opartej na klasteryzacji K-Means).

W metodzie \texttt{generate\_local\_routes}, przesyłki są grupowane według strefy docelowego Paczkomatu:

\begin{lstlisting}[language=Python, caption=Grupowanie przesyłek według stref]
packages_by_zone = defaultdict(list)
for pkg in local_packages:
    zone = pkg.destination_postmat.zone
    if zone:
        packages_by_zone[zone.id].append(pkg)
\end{lstlisting}

\subsection{Teoretyczne ujęcie problemu komiwojażera (TSP)}

Problem komiwojażera (ang. \textit{Travelling Salesman Problem}, TSP) to jedno z fundamentalnych zagadnień optymalizacyjnych w informatyce i badaniach operacyjnych. W swojej klasycznej postaci problem ten formułuje się następująco: mając daną listę miast oraz odległości między każdą parą z nich, należy wyznaczyć najkrótszą trasę, która odwiedza każde miasto dokładnie raz i kończy się w punkcie startowym.

Z matematycznego punktu widzenia, TSP modeluje się przy użyciu grafu ważonego $G = (V, E)$, gdzie $V$ oznacza zbiór wierzchołków (punktów dostaw), a $E$ zbiór krawędzi o określonych wagach (kosztach przejazdu). Rozwiązaniem problemu jest znalezienie cyklu Hamiltona o minimalnej sumie wag krawędzi.

Problem ten należy do klasy problemów NP-trudnych (ang. \textit{NP-hard}). Oznacza to, że nie jest znany algorytm, który potrafiłby znaleźć optymalne rozwiązanie w czasie wielomianowym dla dowolnych danych wejściowych. Wraz ze wzrostem liczby punktów $n$, liczba możliwych permutacji tras rośnie silniowo, co czyni przegląd zupełny (ang. \textit{brute-force}) niepraktycznym dla większych instancji. W zastosowaniach logistycznych, gdzie kluczowy jest czas odpowiedzi systemu, powszechnie stosuje się algorytmy heurystyczne lub metaheurystyczne, które pozwalają znaleźć rozwiązanie suboptymalne (wystarczająco dobre) w akceptowalnym czasie.

\subsection{Implementacja algorytmu TSP i formuła Haversine'a}

Dla każdej strefy generowana jest trasa realizująca problem komiwojażera (ang. travelling salesman problem, TSP). Podobnie jak w logistyce krajowej, zastosowano algorytm najbliższego sąsiada, co zapewnia szybki czas obliczeń przy akceptowalnej jakości tras dla lokalnych dystansów.

Odległości między punktami (Paczkomatami) obliczane są przy użyciu wzoru Haversine'a, która uwzględnia krzywiznę Ziemi, co jest kluczowe przy wykorzystaniu współrzędnych geograficznych (GPS).

\begin{lstlisting}[language=Python, caption=Implementacja wzoru Haversine'a]
def _haversine(self, lat1, lon1, lat2, lon2):
    R = 6371.0  # Promien Ziemi w km
    dlat = math.radians(lat2 - lat1)
    dlon = math.radians(lon2 - lon1)
    a = (math.sin(dlat / 2) ** 2 +
         math.cos(math.radians(lat1)) * math.cos(math.radians(lat2)) *
         math.sin(dlon / 2) ** 2)
    c = 2 * math.atan2(math.sqrt(a), math.sqrt(1 - a))
    return R * c
\end{lstlisting}

Metoda \texttt{\_create\_zone\_route} tworzy sekwencję przystanków (\texttt{RouteStop}), zaczynając i kończąc w magazynie, minimalizując dystans przebyty przez kuriera wewnątrz strefy.

\section{Zarządzanie skrytkami i weryfikacja dostępności}

Kluczowym elementem systemu FasterPost jest dynamiczne zarządzanie dostępnością skrytek w Paczkomatach. Przed zaplanowaniem trasy, system musi zagwarantować, że w docelowym Paczkomacie znajduje się wolna skrytka o odpowiednim rozmiarze.

\subsection{Algorytm alokacji (Allocation Phase)}

Metoda \texttt{\_allocate\_stashes} realizuje logikę dopasowania ("Best Fit" z priorytetem rozmiaru). Algorytm sprawdza dostępne skrytki (\texttt{is\_empty=True}) i rezerwuje je dla konkretnych paczek.

\begin{lstlisting}[language=Python, caption=Logika dopasowania rozmiaru paczki do skrytki]
if p.size == 'small':
    if small_stashes: selected_stash = small_stashes.pop()
    elif medium_stashes: selected_stash = medium_stashes.pop()
    elif large_stashes: selected_stash = large_stashes.pop()
elif p.size == 'medium':
    if medium_stashes: selected_stash = medium_stashes.pop()
    elif large_stashes: selected_stash = large_stashes.pop()
\end{lstlisting}

Jeśli alokacja się powiedzie, skrytka otrzymuje tymczasową rezerwację, która jest finalizowana w momencie utworzenia obiektu \texttt{RoutePackage}. W przypadku braku miejsc (zbiór \texttt{failed\_pkgs}), system automatycznie generuje aktualizację statusu paczki z informacją o opóźnieniu ("Destination locker full"), co pozwala na ponowienie próby w następnym cyklu planowania.
